\chapter{Integers in binary}
\label{chap.integersInBinary}

The book of Section \ref{sec.fasterBook}
keeps the occupied prices of a side as the set bits of one integer,
and finds its best price, and walks its levels, with a few operations on that integer.
This appendix states them as properties of the binary representation, proves them,
and says what each costs.
Most properties are followed by a short program that checks them on many integers,
so that they can be seen as well as proved.

\section{A number as a set}
\label{sec.numberAsSet}
\input{binary_sets}

\section{Bit sequences, two ways}
\label{sec.bitSequences}
\input{binary_sequences}

\section{Two's complement}
\label{sec.twosComplement}
\input{binary_complement}

\section{The integers as bit sequences}
\label{sec.integersAsSequences}
\input{binary_embedding}

\section{Degree, and the lowest set bit}
\label{sec.binaryDegree}
\input{binary_degree}

\section{What each operation costs}
\label{sec.binaryCosts}
\input{binary_costs}

\section{Rank, and the runs of a set}
\label{sec.binaryRuns}
\input{binary_runs}

\section{Where the book uses it}
\label{sec.binaryInTheBook}
\input{binary_in_the_book}
